\documentclass[10pt]{article}
\usepackage[margin=1in]{geometry}
\usepackage{amsmath,amssymb,amsthm,microtype}
\usepackage[colorlinks=true,linkcolor=blue,citecolor=blue,urlcolor=blue]{hyperref}
\newtheorem{theorem}{Theorem}
\newtheorem{lemma}[theorem]{Lemma}
\newtheorem{remark}[theorem]{Remark}
\newcommand{\N}{\mathbb N}
\newcommand{\Z}{\mathbb Z}
\newcommand{\Q}{\mathbb Q}
\newcommand{\C}{\mathbb C}
\newcommand{\R}{\mathbb R}
\newcommand{\vis}{\operatorname{Vis}}
\newcommand{\inv}{\operatorname{Inv}}
\hypersetup{pdftitle={Polynomial visibility has density one},
pdfsubject={An elementary proof and complete Lean formalization}}
\begin{document}
\begin{center}
{\LARGE Polynomial visibility has density one\par}
\vspace{0.3em}
{\large An elementary bounded-gcd argument\par}
\vspace{0.6em}
{\small September 4, 2026}
\end{center}
\begin{abstract}
For every nonzero integral polynomial with at least two distinct complex
roots, almost every positive lattice point is visible along its rational
vertical dilates. The proof has two ingredients: the gcd of a polynomial
value and an independent height is bounded outside a set of arbitrarily
small upper density; and every fixed proper ratio of polynomial values
occurs at a set of arguments of density zero. Both ingredients admit
elementary proofs. No uniform estimate over varying ratios is needed.
\end{abstract}

\section{Statement and reduction}
For $F\in\Z[x]$, call $(a,h)\in\Z_{>0}^2$ \emph{visible} if there is
$t\in\Q$ with $h=tF(a)$ and $a$ is the least positive integer $u$ for
which $tF(u)$ is a positive integer. Put
\[
 V_F(N)=\#\{(a,h)\in[1,N]^2\cap\Z^2:(a,h)\text{ is visible}\}.
\]
\begin{theorem}\label{thm:main}
Let $F\in\Z[x]\setminus\{0\}$ have at least two distinct roots in $\C$.
Then
\[
                    \lim_{N\to\infty}\frac{V_F(N)}{N^2}=1.
\]
\end{theorem}
This is the nonzero version of the Visibility Density Conjecture of
Lobsenz--Phillips \cite{LP}, extending the origin-passing formulation
of Chaubey--Pandey \cite{CP}. Lobsenz--Phillips prove the assertion for
proper powers of polynomials with at least two distinct roots; their
revised paper also records the resolution of the case $F(0)=0$ by
Chaubey, Pandey, and Regavim. The argument below treats all nonzero
polynomials under the stated root
hypothesis. The explicit exclusion of $F=0$ is necessary: otherwise
every complex number is a root, but no positive lattice point is visible.

Replacing $F$ by $-F$ leaves visibility unchanged, by replacing $t$ by
$-t$. We may therefore assume that its leading coefficient is positive.
There is $A_F$ such that, whenever $a>A_F$,
\begin{equation}\label{eq:record}
             F(a)>\max\bigl(0,F(1),\ldots,F(a-1)\bigr).
\end{equation}
Indeed, a nonconstant polynomial with positive leading coefficient is
eventually strictly increasing and tends to infinity. The columns
$a\le A_F$ contain only $O_F(N)$ points.

For $a>A_F$ set $g=\gcd(F(a),h)$. If $(a,h)$ is invisible, its unique
possible scale is $t=h/F(a)$, and some $1\le u<a$ satisfies
$hF(u)/F(a)\in\Z_{>0}$. Consequently
\[
 F(a)\mid hF(u),\qquad 0<F(u)<F(a).
\]
Since $F(a)/g$ and $h/g$ are coprime, $F(a)/g$ divides $F(u)$.
Thus there is an integer $r$ such that
\begin{equation}\label{eq:finite-ratio}
                  F(u)=\frac r g F(a),\qquad 1\le r<g.
\end{equation}
We will show that $g$ can first be bounded at arbitrarily small density
cost, and then handle the resulting finite list of ratios.

\section{The gcd cutoff}
\begin{lemma}\label{lem:tight}
For every nonzero $F\in\Z[x]$ and every $\varepsilon>0$ there is an
integer $M\ge1$ such that, for every $N\ge1$,
\[
 \frac1{N^2}
 \#\{1\le a,h\le N:\gcd(|F(a)|,h)>M\}<\varepsilon.
\]
\end{lemma}
\begin{proof}
Write $d=\deg F$. Choose an integer $K\ge2$ so large that every prime
for which $F$ is the zero polynomial modulo that prime is at most $K$.
This excludes only finitely many primes, since $F$ has a nonzero
coefficient. For any prime $p>K$, the congruence $F(a)\equiv0\pmod p$
has at most $d$ residue classes. Hence
\[
 \#\{1\le a,h\le N:p\mid F(a),\ p\mid h\}
 \le d\left(\frac Np+1\right)\frac Np.
\]
Any prime dividing $h$ is at most $N$, and for $p\le N$ we have
$N/p+1\le2N/p$. Summing over $K<p\le N$ and enlarging the sum to
all integers gives the bound, uniform in $N$,
\begin{equation}\label{eq:large-primes}
 \frac{\#\{(a,h)\in[1,N]^2:\exists p>K,\ p\mid\gcd(|F(a)|,h)\}}{N^2}
 \le 2d\sum_{n>K}\frac1{n^2}\le\frac{2d}{K}.
\end{equation}
For any fixed positive integer $E$, the proportion of heights divisible
by $p^E$ is at most $p^{-E}$. Outside a set of pairs of proportion at most
$\sum_{p\le K}p^{-E}$, none of these small-prime powers divides $h$.
If neither exceptional event occurs, then
\[
       \gcd(|F(a)|,h)\le M:=\prod_{p\le K}p^{E-1}.
\]
Choose $K$ with $2d/K<\varepsilon/2$, and then $E$ with
$\sum_{p\le K}p^{-E}<\varepsilon/2$. Combining the two exceptional
sets proves the claim. This also covers $F(a)=0$,
since $\gcd(0,h)=h>0$ and the same prime-divisor argument applies.
\end{proof}

\section{Fixed ratios are sparse}
\begin{lemma}\label{lem:sparse}
Let $F\in\R[x]$ have degree $d\ge2$, positive leading coefficient,
and at least two distinct complex roots. For each $q\in(0,1)$, the set
\[
 S_q=\{a\in\Z_{>0}:\exists u\in\Z_{>0},\ u<a,\ F(u)=qF(a)\}
\]
has natural density zero.
\end{lemma}
\begin{proof}
Write
\[
 F(x)=c x^d+b x^{d-1}+O_F(x^{d-2}),\quad
 \alpha=q^{1/d}\in(0,1),\quad
 \beta=\frac{b(\alpha-1)}{dc}.
\]
For any fixed $\eta>0$, the polynomials
$F(\alpha x+\beta\pm\eta)-qF(x)$ have degree $d-1$ and leading
coefficients $\pm dc\alpha^{d-1}\eta$. Consequently, for large $a$,
\[
 F(\alpha a+\beta-\eta)<qF(a)<F(\alpha a+\beta+\eta).
\]
Along solutions with $a\to\infty$ we have $u\to\infty$.
Eventual strict monotonicity of $F$ therefore places $u$ between these
two arguments. As $\eta$ is arbitrary, this proves
\begin{equation}\label{eq:asymptote}
                  u=\alpha a+\beta+o_{F,q}(1).
\end{equation}

Suppose first that $\alpha$ is irrational. Fix $L\ge1$ and set
\[
 \delta_L=\min_{1\le j\le L}\operatorname{dist}(j\alpha,\Z)>0.
\]
Beyond a fixed threshold, the error in \eqref{eq:asymptote} has absolute
value less than $\delta_L/3$. Two such solutions with
$1\le a'-a\le L$ would give
\[
 \operatorname{dist}\bigl((a'-a)\alpha,\Z\bigr)
 \le |(u'-u)-\alpha(a'-a)|<2\delta_L/3,
\]
a contradiction. Thus every block of $L+1$ consecutive sufficiently
large integers contains at most one member of $S_q$. Its upper density
is at most $1/(L+1)$; letting $L\to\infty$ proves the assertion.

Suppose next that $\alpha$ is rational. Choose $D\ge1$ such that
$D\alpha\in\Z$. Then $u-\alpha a\in D^{-1}\Z$.
Equation \eqref{eq:asymptote} shows that, if arbitrarily large solutions
exist, necessarily $\beta\in D^{-1}\Z$ and all sufficiently large
solutions satisfy $u=\alpha a+\beta$ exactly. If infinitely many
solutions existed, the polynomial
\[
                    F(\alpha x+\beta)-qF(x)
\]
would have infinitely many zeros, and so would vanish identically.
Let $R$ be the finite nonempty set of complex roots of $F$, and let
$T(z)=\alpha z+\beta$. This identity implies $T(R)\subseteq R$.
As $T$ is injective and $R$ finite, $T(R)=R$. But the diameter satisfies
\[
          \operatorname{diam}(R)=\operatorname{diam}(T(R))
                              =\alpha\operatorname{diam}(R).
\]
Since $0<\alpha<1$, the diameter is zero, contrary to the existence of
two distinct roots. In this case $S_q$ is in fact finite.
\end{proof}

\section{Completion of the proof}
\begin{proof}[Proof of Theorem~\ref{thm:main}]
Fix $\varepsilon>0$ and choose $M$ from Lemma~\ref{lem:tight}.
Let
\[
       S_M=\bigcup_{2\le g\le M}\ \bigcup_{1\le r<g}S_{r/g}.
\]
This is a finite union of sets of density zero by Lemma~\ref{lem:sparse},
so $\#(S_M\cap[1,N])=o_{F,M}(N)$.
By \eqref{eq:finite-ratio}, every invisible point in $[1,N]^2$ either
has $a\le A_F$, has $\gcd(F(a),h)>M$, or has $a\in S_M$.
It follows that
\[
 \limsup_{N\to\infty}\frac{N^2-V_F(N)}{N^2}
 \le 0+\varepsilon+
       \lim_{N\to\infty}\frac{N\,\#(S_M\cap[1,N])}{N^2}
 =\varepsilon.
\]
Since $\varepsilon$ is arbitrary and the left side is nonnegative,
$N^2-V_F(N)=o(N^2)$.
\end{proof}

\begin{remark}
The order of limits is essential: first fix an error tolerance and its
gcd cutoff $M$, then let $N\to\infty$. The proof requires no uniform
control of $S_{r/g}$ as $g$ grows. In particular it does not assert an
unproved average-gcd bound for two polynomial values.
\end{remark}

\paragraph{Formal verification.}
The accompanying Lean~4.28.0 project proves the integer-polynomial theorem
as \texttt{PolynomialVisibility.visibility\_density\_one}, using
mathlib~v4.28.0. The complete proof has no admitted steps or additional
axioms; its dependencies are only the standard logical foundations
\texttt{propext}, \texttt{Classical.choice}, and \texttt{Quot.sound}.

\begin{thebibliography}{9}
\bibitem{CP} S. Chaubey and A. K. Pandey,
\emph{On the density of visible lattice points along polynomials},
2021, \href{https://arxiv.org/abs/2109.08431}{arXiv:2109.08431}.
\bibitem{LP} A. Lobsenz and T. Phillips,
\emph{Lattice point visibility along powers of polynomials},
2026, \href{https://arxiv.org/abs/2604.23050v3}{arXiv:2604.23050v3}.
\end{thebibliography}
\end{document}
